\documentclass[acmsmall,nonacm,screen]{acmart}

\settopmatter{printacmref=false}
\setcopyright{none}
\renewcommand\footnotetextcopyrightpermission[1]{}

\citestyle{acmauthoryear}

\usepackage{mathpartir}   % \mathpar / \inferrule for inductive relation rules
\usepackage{placeins}

\newcommand{\romem}{$Rome^{M}$}
\newcommand{\romam}{$Roma^{M}$}
\newcommand{\reggiom}{$Reggio^{M}$}
\newcommand{\inN}{\!\in\!}
\newcommand{\tick}{\checkmark}

\begin{document}

\title{Verona through Rome}

\author{Hoang Nguyen}
% \affiliation{%
%   \institution{Imperial College London}
%   \city{London}
%   \country{United Kingdom}
% }
\email{hoang4j@gmail.com}

\begin{abstract}
\end{abstract}

\maketitle

\section{Introduction}

We present \romam, a more complete, fully formalized version of \romem \cite{nguyen_verona_2026}.

\section{Model}

The configuration remains unchanged from the report. We make one change to the semantics, specifically within \texttt{makeRegion}:

\begin{mathpar}
  \inferrule* [lab=make-region]
    {F = (rid, bvar, omap, vmap, fid) \\ \langle (S :: F), H \rangle.freshObjectId = oid' \\ \mathbf{\langle (S :: F), H \rangle.freshRegionId = rid'} \\ R' = (oid', [oid' \mapsto \emptyset], Closed) \\ F' = (rid, bvar, omap, vmap[\texttt{x} \mapsto \mathsf{RId}(rid')], fid)}
    {\langle (S :: F), H \rangle \xrightarrow{\texttt{makeRegion x}} \langle S :: F', H[rid \mapsto R'] \rangle}
\end{mathpar}

In the report, the \texttt{regionId} would coincide with the \texttt{objectId} of the bridge object, but here we've chosen to completely decouple them. The original reasoning behind this was to future-proof for deallocation, but in writing this we found out that the argument didn't hold, and it would've been fine to leave \texttt{makeRegion} as before. Luckily, this doesn't change much, it only means that we need to prove regions are unique separately from objects, as opposed to using a combined argument before.

\section{Invariants}
\label{sec:invariants}

\begin{definition}[L1]
No two objects in the configuration share the same identifier.
\[\langle S, H \rangle \vdash L1 \triangleq \langle S, H \rangle.objectIds \text{ is duplicate-free.}\]
\end{definition}

\begin{definition}[L2]
Every frame on the stack refers to an open region.
\[\langle S, H \rangle \vdash L2 \triangleq \forall frame \inN S : \exists region : [H(frame.regionId) = region \wedge region.status = Open]\]
\end{definition}

\begin{definition}[H1]
A region always contains its own bridge object.
\[\langle S, H \rangle \vdash H1 \triangleq \forall region \inN H.regions : [region.bridgeObjectId \inN dom(region.objMap)]\]
\end{definition}

\begin{definition}[H2]
There is at most one reference to a region from anywhere in the configuration.
\[\langle S, H \rangle \vdash H2 \triangleq \forall rid : [S.refs.count(\mathsf{RId}(rid)) + H.refs.count(\mathsf{RId}(rid)) \le 1]\]
\end{definition}

\begin{definition}[H3]
A reference stored inside a region resolves back into that same region.
\begin{align*}
  \langle S, H \rangle \vdash H3 \triangleq{} & \forall rid, oid, region : [H(rid) = region \wedge \mathsf{OId}(oid) \inN region.refs \implies{} \\
  & \mathsf{OId}(oid).loc(\langle S, H \rangle) = \mathsf{Rgn}(rid)]
\end{align*}
\end{definition}

\begin{definition}[S1]
No two frames on the stack share a region.
\[\langle S, H \rangle \vdash S1 \triangleq \forall f_1, f_2 \inN S : [f_1.regionId = f_2.regionId \implies f_1 = f_2]\]
\end{definition}

\begin{definition}[S2]
A reference from a frame to a temporary object resolves to that frame or an earlier one.
\begin{align*}
  \langle S, H \rangle \vdash S2 \triangleq{} & \forall frame \inN S, \mathsf{OId}(oid) \inN frame.refs, fid' :\ [\mathsf{OId}(oid).loc(\langle S, H \rangle) = \mathsf{Stk}(fid') \implies{} \\
  & fid' \le frame.fid]
\end{align*}
\end{definition}

\begin{definition}[S3]
A reference from a frame to a region object resolves to that frame's region or the region of an earlier frame.
\begin{align*}
  \langle S, H \rangle \vdash S3 \triangleq{} & \forall frame \inN S, \mathsf{OId}(oid) \inN frame.refs, rid' :\ [\mathsf{OId}(oid).loc(\langle S, H \rangle) = \mathsf{Rgn}(rid') \implies{} \\
  & \exists frame' \inN S : [frame'.regionId = rid' \wedge frame'.fid \le frame.fid]]
\end{align*}
\end{definition}

We add two new invariants:

\begin{definition}[HS1]
  Every reference to an object must be to an object that exists in the configuration.
  \[\langle S, H \rangle \vdash HS1 \triangleq \forall oid : [\mathsf{OId}(oid) \inN \langle S, H \rangle.refs \implies oid \inN \langle S, H \rangle.objectIds]\]
\end{definition}

\begin{definition}[HS2]
  Every reference to a region must be to a region that exists in the configuration.
  \[\langle S, H \rangle \vdash HS2 \triangleq \forall rid : [\mathsf{RId}(rid) \inN \langle S, H \rangle.refs \implies rid \inN dom(H)]\]
\end{definition}

These came up in the process of proving \textit{S2} and \textit{S3} for \texttt{makeObjStack}, \texttt{makeObjRegion}, and \texttt{makeRegion}. This stems from the fact that our \texttt{freshObjectId} and \texttt{freshRegionId} functions only take into account the \texttt{Ids} used as keys when creating a new \texttt{Id}. Thus, without this invariant, it is technically possible\footnote{Although it isn't possible within the context of an operational semantics starting with the starting configuration, which is why we were able to prove it without altering the operational semantics.} for a configuration to contain a dangling reference to an \texttt{objectId} or \texttt{regionId} that does not exist within the configuration, but is then inserted via one of the three \texttt{make} operations.

\begin{definition}[ValidConfig]
\label{def:validconfig}
\begin{mathpar}
  \inferrule* [lab=Valid]
    {\langle S,H\rangle \vdash L1 \\ \langle S,H\rangle \vdash L2 \\ \langle S,H\rangle \vdash H1 \\ \langle S,H\rangle \vdash H2 \\ \langle S,H\rangle \vdash H3 \\ \langle S,H\rangle \vdash S1 \\ \langle S,H\rangle \vdash S2 \\ \langle S,H\rangle \vdash S3 \\ \langle S,H\rangle \vdash HS1 \\ \langle S,H\rangle \vdash HS2}
    {\langle S,H\rangle \vdash \mathit{ValidConfig}}
\end{mathpar}
\end{definition}

\texttt{RuntimeConfig.start} (one root region, one root frame) satisfies \textit{ValidConfig}, and every one of the nine mutation operations of \S\ref{sec:opsem} preserves it; consequently every reachable configuration of the operational semantics is a \textit{ValidConfig}.

\section{Reachability}
\label{sec:reachability}
We introduce the notion of reachability. It is the same term as in the original report, but the definition here covers a wider spectrum.

\begin{definition}[Stack-reachable]
  An object or region is stack-reachable if it is frame-reachable from some frame in the stack.
  \begin{mathpar}
    \inferrule*
      {\langle S, H \rangle \vdash FReachable(\_, a)}
      {\langle S, H \rangle \vdash SReachable(a)}
  \end{mathpar}
\end{definition}

\begin{definition}[Frame-reachable]
  An object or region is frame-reachable if there is a path from a frame to said object or region.
  \begin{mathpar}
    \inferrule*[lab=var]
      {F = (\_, \_, \_, vmap, fid) \\ vmap(\_) = a}
      {\langle S :: F, H \rangle \vdash FReachable(fid, a)}

    \inferrule*[lab=bridge]
      {F = (rid, \_, \_, \_, fid) \\ H(rid) = (boid, \_, \_)}
      {\langle S :: F, H \rangle \vdash FReachable(fid, \mathsf{OId}(boid))}

    \inferrule*[lab=step]
      {\langle S, H \rangle \vdash FReachable(fid, a) \\ \langle S, H \rangle \vdash a \to b}
      {\langle S, H \rangle \vdash FReachable(fid, b)}
  \end{mathpar}
\end{definition}

\begin{definition}[Region-reachable]
  An object or region is region-reachable if there is a path from the bridge object to said object or region.
  \begin{mathpar}
    \inferrule*[lab=bridge]
      {H(rid) = (boid, \_, \_)}
      {\langle S, H \rangle \vdash RReachable(rid, \mathsf{OId}(boid))}

    \inferrule*[lab=step]
      {\langle S, H \rangle \vdash RReachable(rid, a) \\ \langle S, H \rangle \vdash a \to b}
      {\langle S, H \rangle \vdash RReachable(rid, b)}
  \end{mathpar}
\end{definition}

\begin{definition}[Reachable-step]
  A one-step reachable relation between objects and regions. An object or region is one-step reachable from another object if that object contains a reference to said object or region. An object is one-step reachable from a region if the region is closed and the object is the region's bridge object.
  \begin{mathpar}
    \inferrule*[lab=obj-step]
      {\mathsf{OId}(oid).objAt(\langle S, H \rangle) = obj \\ b \inN obj.refs }
      {\langle S, H \rangle \vdash \mathsf{OId}(oid) \to b}

    \inferrule*[lab=region-step]
      {H(rid) = (boid, \_, Closed)}
      {\langle S, H \rangle \vdash \mathsf{RId}(rid) \to \mathsf{OId}(boid)}
  \end{mathpar}
\end{definition}

\begin{theorem}
  Region-reachability is the transitive closure of the reachable-step from the region's bridge object.
  \[\langle S, H \rangle \vdash RReachable(rid, a) \iff \exists boid : [H(rid) = (boid, \_, \_) \wedge \mathsf{OId}(boid) \xrightarrow{*} a]\]
\end{theorem}

\begin{theorem}
  Frame-reachability is the transitive closure of the reachable-step from the frame.
  \[\langle S, H \rangle \vdash FReachable(fid, a) \iff \exists start : [\langle S, H \rangle \vdash FRoot(fid, start) \wedge start \xrightarrow{*} a]\]
\end{theorem}

\begin{definition}[Frame-root]
  The root coming from a frame can be a reference stored by a variable or the bridge object of that frame's region.
  \begin{mathpar}
    \inferrule*[lab=var]
      {F = (\_, \_, \_, vmap, fid) \\ vmap(\_) = a}
      {\langle S :: F, H \rangle \vdash FRoot(fid, a)}

    \inferrule*[lab=bridge]
      {F = (rid, \_, \_, \_, fid) \\ H(rid) = (boid, \_, \_)}
      {\langle S :: F, H \rangle \vdash FRoot(fid, \mathsf{OId}(boid))}
  \end{mathpar}
\end{definition}

\subsection{Suspended regions}
\label{sec:suspended}

\begin{definition}[FrameReachableAtLaterFrame]
\label{def:fralf}
A configuration satisfies \textit{FrameReachableAtLaterFrame} if, whenever a frame $F$ is earlier on the stack than a frame $F'$, any object located in $F$'s region that is frame-reachable from $F'$ is already frame-reachable from $F$ itself.
\begin{align*}
  \mathit{FRALF}(\langle S, H \rangle) \iff{} & \forall F, F' \inN S.\ F.fid < F'.fid \implies{} \\
  & \forall oid.\ \mathsf{OId}(oid).loc(\langle S, H \rangle) = \mathsf{Rgn}(F.regionId) \implies{} \\
  & \langle S, H \rangle \vdash FReachable(F'.fid, \mathsf{OId}(oid)) \implies \langle S, H \rangle \vdash FReachable(F.fid, \mathsf{OId}(oid))
\end{align*}
\end{definition}

\begin{theorem}[FRALF is preserved]
\label{thm:fralf-step}
Let $cmd$ be any of the nine operations of \S\ref{sec:opsem}. If $\langle S, H \rangle$ is a \textit{ValidConfig} satisfying \textit{FRALF}, and $\langle S, H \rangle \xrightarrow{cmd} \langle S', H' \rangle$, then $\langle S', H' \rangle$ also satisfies \textit{FRALF}.
\end{theorem}

\begin{theorem}[Suspended-region stack-reachability is invariant]
\label{thm:suspended-invariant}
Let $cmd$ be any of the nine operations of \S\ref{sec:opsem}, and let $\langle S, H \rangle$ be a \textit{ValidConfig} satisfying \textit{FRALF} with $\langle S, H \rangle \xrightarrow{cmd} \langle S', H' \rangle$. For any suspended frame $F \inN S$ and any object $oid$ located in $F$'s region,
\[
  \langle S, H \rangle \vdash SReachable(\mathsf{OId}(oid)) \iff \langle S', H' \rangle \vdash SReachable(\mathsf{OId}(oid))
\]
\end{theorem}

\section{Related Work}

\section{Conclusion}

\appendix
\section{Configuration}
{\setlength{\arraycolsep}{2pt}
\[
\begin{array}{rcl}
  RegionId, ObjectId, Index & \subseteq & Int \\
  VarName, FieldName & \subseteq & String \\
  &&\\
  RuntimeConfiguration & \triangleq & Heap \times Stack\\
  &&\\
  Heap & \triangleq & RegionId \rightharpoonup Region\\
  Region & \triangleq & ObjectId\times ObjMap \times Status\\
  Status & \triangleq & Open \ | \ Closed\\
  &&\\
  Stack & \triangleq & List \ Frame\\
  Frame & \triangleq & RegionId \times VarName \times ObjMap \times VarMap \times Index\\
  &&\\
  ObjMap & \triangleq & ObjectId \rightharpoonup FieldMap\\
  &&\\
  Object & \triangleq & FieldMap\\
  FieldMap & \triangleq & FieldName \rightharpoonup Reference\\
  VarMap & \triangleq & VarName \rightharpoonup Reference\\
  &&\\
  Reference & \triangleq & \mathsf{RId}(RegionId) \ | \ \mathsf{OId}(ObjectId)
\end{array}
\]}

\section{Operational Semantics}
\label{sec:opsem}
\begin{mathpar}
  \inferrule* [lab=var-asgn-stack]
    {F = (rid, bvar, omap, vmap, fid) \\ \texttt{x} \ne bvar \\ resolveV(\texttt{x}, \langle (S :: F), H \rangle) = \epsilon \\ resolveFA(\texttt{y.f}, \langle (S :: F), H \rangle) = \mathsf{OId}(oid)}
    {\langle (S :: F), H \rangle \xrightarrow{\texttt{x := y.f}} \langle S :: (rid, bvar, omap, vmap[\texttt{x} \mapsto \mathsf{OID}(oid)], fid), H \rangle}
  \\
  \inferrule* [lab=var-asgn-bridge]
    {F = (rid, \texttt{x}, \_, \_, \_) \\ resolveFA(\texttt{y.f}, \langle (S :: F), H \rangle) = \mathsf{OId}(oid) \\ \mathsf{OId}(oid).loc(\langle (S :: F), H \rangle) = \mathsf{Rgn}(rid) \\ H(rid) = (boid, omap, status)}
    {\langle (S :: F), H \rangle \xrightarrow{\texttt{x := y.f}} \langle (S :: F), H[rid \mapsto (oid, omap, status)] \rangle}
  \\
\end{mathpar}
\begin{mathpar}
  \inferrule* [lab=field-asgn-stack]
    {F = (rid, bvar, omap, vmap, fid) \\ resolveV(\texttt{x}, \langle (S :: F), H \rangle) = \mathsf{OId}(oid) \\ \mathsf{OId}(oid).loc(\langle (S :: F), H \rangle) = \mathsf{Stk}(fid) \\ resolveV(\texttt{y}, \langle (S :: F), H \rangle) = \mathsf{OId}(oid') \\ omap(oid) = obj}
    {\langle (S :: F), H \rangle \xrightarrow{\texttt{x.f := y}} \langle (S :: (rid, bvar, omap[oid \mapsto obj[\texttt{f} \mapsto OId(oid')]], vmap, fid)), H \rangle}
  \\
  \inferrule* [lab=field-asgn-region]
    {F = (rid, \_, \_, \_, \_) \\ resolveV(\texttt{x}, \langle (S :: F), H \rangle) = \mathsf{OId}(oid) \\ \mathsf{OId}(oid).loc(\langle (S :: F), H \rangle) = \mathsf{Rgn}(rid) \\ resolveV(\texttt{y}, \langle (S :: F), H \rangle) = \mathsf{OId}(oid') \\ \mathsf{OId}(oid').loc(\langle (S :: F), H \rangle) = \mathsf{Rgn}(rid) \\ H(rid) = (boid, omap, Open) \\ omap(oid) = obj }
    {\langle (S :: F), H \rangle \xrightarrow{\texttt{x.f := y}} \langle (S :: F), H[rid \mapsto (boid, omap[oid \mapsto obj[\texttt{f} \mapsto \mathsf{OId}(oid')]], Open)] \rangle}
  \\
\end{mathpar}
\begin{mathpar}
  \inferrule* [lab=swap-stack]
    {F = (rid, bvar, omap, vmap, fid) \\ \texttt{x} \ne bvar \\ vmap(\texttt{x}) = \texttt{x}ref \\ resolveV(\texttt{y}, \langle (S :: F), H \rangle) = \mathsf{OId}(\texttt{y}oid) \\ resolveFA(\texttt{y.f}, \langle (S :: F), H \rangle) = \texttt{yf}ref \\ \mathsf{OId}(\texttt{y}oid).loc(\langle (S :: F), H \rangle) = \mathsf{Stk}(fid) \\ omap(\texttt{y}oid) = obj \\ omap' = omap[\texttt{y}oid \mapsto obj[\texttt{f} \mapsto \texttt{x}ref]] \\ vmap' = vmap[\texttt{x} \mapsto \texttt{yf}ref] \\ F' = (rid, bvar, omap', vmap')}
    {\langle (S :: F), H \rangle \xrightarrow{\texttt{x.f := y}} \langle (S :: F'), H \rangle}
  \\
  \inferrule* [lab=swap-region-object]
    {F = (rid, bvar, fomap, vmap, fid) \\ \texttt{x} \ne bvar \\ vmap(\texttt{x}) = \mathsf{OId}(\texttt{x}oid) \\ resolveV(\texttt{y}, \langle (S :: F), H \rangle) = \mathsf{OId}(\texttt{y}oid) \\ resolveFA(\texttt{y.f}, \langle (S :: F), H \rangle) = \texttt{yf}ref \\ \{\mathsf{OId}(\texttt{x}oid), \mathsf{OId}(\texttt{y}oid)\}.loc(\langle (S :: F), H \rangle) = \mathsf{Rgn}(rid) \\ H(rid) = (boid, romap, Open) \\ romap(\texttt{y}oid) = obj \\ obj' = obj[\texttt{f} \mapsto \mathsf{OId}(\texttt{x}oid)] \\ R' = (boid, romap[\texttt{y}oid \mapsto obj'], Open) \\ F' = (rid, bvar, fomap, vmap[\texttt{x} \mapsto \texttt{yf}ref])}
    {\langle (S :: F), H \rangle \xrightarrow{\texttt{x.f := y}} \langle (S :: F'), H[rid \mapsto R'] \rangle}
  \\
  \inferrule* [lab=swap-region-region]
    {F = (rid, bvar, fomap, vmap, fid) \\ \texttt{x} \ne bvar \\ vmap(\texttt{x}) = \mathsf{RId}(\texttt{x}rid) \\ resolveV(\texttt{y}, \langle (S :: F), H \rangle) = \mathsf{OId}(\texttt{y}oid) \\ resolveFA(\texttt{y.f}, \langle (S :: F), H \rangle) = \texttt{yf}ref \\ \mathsf{OId}(\texttt{y}oid).loc(\langle (S :: F), H \rangle) = \mathsf{Rgn}(rid) \\ H(xrid) = (\_, \_, Closed) \\ H(rid) = (boid, romap, Open) \\ romap(\texttt{y}oid) = obj \\ obj' = obj[\texttt{f} \mapsto \mathsf{RId}(\texttt{x}rid)] \\ R' = (boid, romap[\texttt{y}oid \mapsto obj'], Open) \\ F' = (rid, bvar, fomap, vmap[\texttt{x} \mapsto \texttt{yf}ref])}
    {\langle (S :: F), H \rangle \xrightarrow{\texttt{x.f := y}} \langle (S :: F'), H[rid \mapsto R'] \rangle}
  \\
  \inferrule* [lab=swap-region-bridge]
    {F = (rid, bvar, fomap, vmap, fid) \\ \texttt{x} = bvar \\ resolveV(\texttt{y}, \langle (S :: F), H \rangle) = \mathsf{OId}(\texttt{y}oid) \\ resolveFA(\texttt{y.f}, \langle (S :: F), H \rangle) = \mathsf{OId}(\texttt{yf}oid) \\ \{\mathsf{OId}(\texttt{y}oid),\mathsf{OId}(\texttt{yf}oid)\}.loc(\langle (S :: F), H \rangle) = \mathsf{Rgn}(rid) \\ H(rid) = (boid, romap, Open) \\ romap(\texttt{y}oid) = obj \\ obj' = obj[\texttt{f} \mapsto \mathsf{OId}(boid)] \\ R' = (\texttt{yf}oid, romap[\texttt{y}oid \mapsto obj'], Open)}
    {\langle (S :: F), H \rangle \xrightarrow{\texttt{x.f := y}} \langle (S :: F), H[rid \mapsto R'] \rangle}
  \\
\end{mathpar}
\begin{mathpar}
  \inferrule* [lab=enter]
    {F = (rid, bvar, fomap, vmap, fid) \\ resolveFA(\texttt{x.f}, \langle S :: F, H \rangle) = \mathsf{RId}(rid') \\ H(rid') = (boid, romap, Closed) \\ F' = (rid', \texttt{a}, \emptyset, \emptyset, fid + 1)}
    {\langle (S :: F), H \rangle \xrightarrow{\texttt{enter x.f a}} \langle (S :: F :: F'), H[rid' \mapsto (boid, romap, Open)] \rangle}
  \\
  \inferrule* [lab=exit]
    {F = (rid, \_, \_, \_, \_) \\ H(rid) = (boid, omap, Open)}
    {\langle (S :: F' :: F), H \rangle \xrightarrow{\texttt{exit}} \langle S :: F', H[rid \mapsto (boid, omap, Closed)] \rangle}
  \\
\end{mathpar}
\begin{mathpar}
  \inferrule* [lab=make-obj-stack]
    {F = (rid, bvar, omap, vmap, fid) \\ \langle (S :: F), H \rangle.freshObjectId = oid}
    {\langle (S :: F), H \rangle \xrightarrow{\texttt{makeObjStack x}} \langle S :: (rid, bvar, omap[oid \mapsto \emptyset], vmap[\texttt{x} \mapsto \mathsf{OId}(oid)], fid), H \rangle}
  \\
  \inferrule* [lab=make-obj-region]
    {F = (rid, bvar, fomap, vmap, fid) \\ \langle (S :: F), H \rangle.freshObjectId = oid \\ H(rid) = (boid, romap, Open) \\ R = (boid, romap[oid \mapsto \emptyset], Open)}
    {\langle (S :: F), H \rangle \xrightarrow{\texttt{makeObjRegion x}} \langle S :: (rid, bvar, fomap, vmap[\texttt{x} \mapsto \mathsf{OId}(oid)], fid), H[rid \mapsto R] \rangle}
  \\
  \inferrule* [lab=make-region]
    {F = (rid, bvar, omap, vmap, fid) \\ \langle (S :: F), H \rangle.freshObjectId = oid' \\ \langle (S :: F), H \rangle.freshRegionId = rid' \\ R' = (oid', [oid' \mapsto \emptyset], Closed) \\ F' = (rid, bvar, omap, vmap[\texttt{x} \mapsto \mathsf{RId}(rid')], fid)}
    {\langle (S :: F), H \rangle \xrightarrow{\texttt{makeRegion x}} \langle S :: F', H[rid \mapsto R'] \rangle}
  \\
\end{mathpar}
\begin{mathpar}
  \inferrule* [lab=merge]
    {F = (rid, bvar, omap, vmap, fid) \\ vmap(\texttt{x}) = \mathsf{RId}(rid') \\ H(rid) = (boid, romap, Open) \\ H(rid') = (boid', romap', Closed) \\ F' = (rid, bvar, omap, vmap[\texttt{x} \mapsto boid'], fid) \\ R = (boid, romap \cup romap', Open)}
    {\langle (S :: F), H \rangle \xrightarrow{\texttt{merge x}} \langle S :: F', H[rid' \mapsto \epsilon][rid \mapsto R] \rangle}
\end{mathpar}

% --- Pattern 1: a set-builder / cardinality-style invariant over a config ---
\begin{definition}[Name]
\label{def:example-invariant}
A configuration $\sigma$ satisfies \textnormal{P} if
\[
  \forall x.\quad
  \bigl|\{\, y \mid y \in X(\sigma),\ y = x \,\}\bigr| \le 1.
\]
\end{definition}

% --- Pattern 2: a universally-quantified implication invariant ---
\begin{definition}[Name]
\label{def:example-invariant-2}
\[
  \forall x \in X(\sigma).\ \exists y.\ P(x, y) \land Q(y).
\]
\end{definition}

% --- Pattern 3: an inductive step relation, via mathpartir's inferrule ---
\begin{definition}[Step relation]
\label{def:example-step}
\begin{mathpar}
\inferrule
  {P(\sigma, a) \\ Q(a, f, b)}
  {\sigma \vdash a \xrightarrow{f} b}
\end{mathpar}
\end{definition}

% --- Pattern 4: closure of a relation, defined via iff ---
\begin{definition}[Reachable]
\label{def:example-reachable}
\[
  \mathit{Reachable}(\sigma, i, b)
  \iff
  \exists a.\ \mathit{Root}(\sigma, i, a) \land \sigma \vdash a \xrightarrow{*} b.
\]
\end{definition}

% --- Pattern 7: a two-branch step relation, one rule per case, with a side
% condition gating the second branch (cf. an OId-step vs.\ an RId-step that
% only fires when some guard holds). Two \inferrule's inside one \mathpar
% render side by side; \\ inside the premises stacks them vertically. ---
\begin{definition}[Two-branch step relation]
\label{def:example-step-2branch}
\begin{mathpar}
\inferrule[Step-Object]
  {P(\sigma, a) = c \\ c.f = b}
  {\sigma \vdash a \xrightarrow{r} b}

\inferrule[Step-Region]
  {Q(\sigma, a) = c \\ \mathit{guard}(c) \\ b = \mathit{target}(c)}
  {\sigma \vdash a \xrightarrow{r} b}
\end{mathpar}
\end{definition}

% --- Pattern 8: an inductive relation presented as its own constructors
% (base case(s) + closure case), rather than collapsed straight to the
% \exists\,\mathrm{Root} \land \to^* form of Pattern 4. Prefer this when the
% distinct roots/base cases are worth keeping visible; state the Pattern-4
% form afterward as a derived remark/lemma if useful. ---
\begin{definition}[Inductively-closed reachability]
\label{def:example-inductive-reachable}
\begin{mathpar}
\inferrule[Base-A]
  {x \in X(\sigma) \\ P(x) = a}
  {\mathit{Reach}(\sigma, i, a)}

\inferrule[Base-B]
  {x \in X(\sigma) \\ Q(x) = a}
  {\mathit{Reach}(\sigma, i, a)}

\inferrule[Step]
  {\sigma \vdash a \xrightarrow{f} b \\ \mathit{Reach}(\sigma, i, a)}
  {\mathit{Reach}(\sigma, i, b)}
\end{mathpar}
\end{definition}

% --- Pattern 5: a theorem statement, informal-prose style ---
\begin{theorem}[Name]
\label{thm:example}
Statement.
\end{theorem}

% --- Pattern 6: a theorem with explicit hypotheses/conclusion via cases-style layout ---
\begin{theorem}[Name]
\label{thm:example-2}
Let $\sigma$ be such that $P(\sigma)$. Then
\[
  Q(\sigma) \implies R(\sigma).
\]
\end{theorem}


\newpage

\bibliographystyle{ACM-Reference-Format}
\bibliography{refs}

\end{document}
